% Generated draft: requires mathematical and source review.
\documentclass[11pt]{article}
\usepackage[margin=25mm]{geometry}
\usepackage{fontspec}
\setmainfont{Latin Modern Roman}
\usepackage{amsmath,amssymb,mathtools,unicode-math}
\setmathfont{Latin Modern Math}
\usepackage{enumitem,needspace,xurl,hyperref}
\setlength{\parindent}{0pt}
\setlength{\parskip}{5pt}
\setlength{\emergencystretch}{3em}
\newcommand{\N}{\mathbb{N}}
\newcommand{\Z}{\mathbb{Z}}
\newcommand{\Q}{\mathbb{Q}}
\newcommand{\R}{\mathbb{R}}
\newcommand{\C}{\mathbb{C}}
\newcommand{\F}{\mathbb{F}}
\newcommand{\A}{\mathbb{A}}
\newcommand{\PP}{\mathbb{P}}
\newcommand{\E}{\mathbb{E}}
\newcommand{\eps}{\varepsilon}
\newcommand{\dd}{\,\mathrm{d}}
\begin{document}
\section*{UnsolvedMath ID: 30001173}
\textbf{Source problem code:} catalogue entry
\subsection{Definitions and mathematical statement}
A category consists of objects, morphisms with specified domains and codomains, an identity morphism for each object, and composition of composable morphisms, satisfying associativity and the identity laws. Category theory studies categories and their relationships. This terminology alone does not define a measurable notion of its use in mathematical research.

The supplied record poses the following methodological and historical questions about category theory in contemporary mathematics:
\begin{enumerate}
\item Is it true that category theory is not widely used?
\item If it is not widely used, why is this so?
\item Can one identify branches of mathematics that use category theory and branches that do not?
\item For what kinds of mathematical problems is category theory used?
\end{enumerate}
The premise in the second question is conditional, not an asserted fact. The source supplies no definition of ``widely used,'' no criterion for whether a branch uses category theory, no fixed classification of branches or problems, and no specified population of mathematical work over which usage is to be assessed. The temporal scope of ``contemporary'' is also unspecified; the cited workshop dates from 2009. Consequently, these questions cannot be converted into a uniquely specified mathematical proposition with precise domains and quantifiers from the supplied evidence alone. They are retained as open-ended research questions, without introducing a quantitative usage threshold or asserting their premise.
\subsection{Short English statement}
The record asks how widely category theory is used, why its use might be limited, which mathematical branches use it, and what problems it addresses. These are methodological and historical questions. A precise mathematical proposition would require additional definitions of usage, prevalence, scope, and time period.
\subsection{Sources}
\begin{itemize}
\item Supplied local UnsolvedMath registry record: permanent ID \(30001173\), title \emph{Mathematical Uses of Category Theory}, source problem code \texttt{catalogue entry}, in the collection \emph{Oberwolfach Reports Open Problems}. No individual proposer is identified; the supplied proposed year is 2009.
\item Supplied report citation: \emph{Mini-Workshop: Category Theory and Related Fields: History and Prospects} (2009), \emph{Oberwolfach Reports}, pp.~463--492, DOI: \url{https://doi.org/10.4171/owr/2009/08}. The supplied background describes Oberwolfach Reports as documenting mathematical workshops and their open-problem discussions.
\item Supplied literature source, labelled ``OWR Category Theory and Related Fields'': \url{https://doi.org/10.4171/owr/2009/08}. The literature assessment dated \texttt{2026-08-21} labels the record ``open'' and describes the questions as methodological and historical, remaining open-ended rather than constituting a resolved theorem. This is supplied dated curation derived from the OWR record, not an independently verified assessment; the supplied background requests independent verification of time-sensitive status and cited claims.
\item The registry labels the statement ``exact'' and reports that its mathematical content was checked against the official Oberwolfach report, with notation, context, and question form normalized without changing the source-level claim. This verification description is supplied provenance and has not been independently checked here.
\end{itemize}
\end{document}
